\documentclass[aspectratio=169]{beamer}
\usetheme{metropolis}
\usepackage{graphicx}
\usepackage{tikz}
\usepackage{amsmath}
\usepackage{listings}
\usepackage{xcolor}
\usepackage{booktabs}
\usepackage{emoji}

% Code listing setup
\lstset{
  basicstyle=\ttfamily\small,
  keywordstyle=\color{blue},
  commentstyle=\color{gray},
  stringstyle=\color{red},
  showstringspaces=false,
  breaklines=true,
  frame=single,
  language=Python
}

\title{Cognitive Models of Semantic Representation}
\subtitle{Lecture 11: How Humans Represent Meaning}
\author{LLM Course - Week 4}
\date{Winter 2026}

\begin{document}

% ============================================================
% TITLE SLIDE
% ============================================================
\begin{frame}
\titlepage
\end{frame}

% ============================================================
% OVERVIEW
% ============================================================
\begin{frame}{Today's Lecture 📋}
\large
\begin{enumerate}
    \item 🧠 \textbf{Human vs. Computational Semantics}
    \item 🔬 \textbf{Cognitive Theories of Meaning}
    \item 🌐 \textbf{Embodied \& Grounded Cognition}
    \item 📊 \textbf{Semantic Similarity: What Does It Mean?}
    \item 🔍 \textbf{Empirical Evidence from Cognitive Science}
    \item 🤖 \textbf{Bridging the Gap: Models Meet Minds}
\end{enumerate}

\vspace{1em}
\textit{Goal: Understand what computational models are really learning}
\end{frame}

% ============================================================
% FUNDAMENTAL QUESTION
% ============================================================
\begin{frame}{The Fundamental Question 🤔}
\begin{center}
\LARGE
\textbf{What does "meaning" actually mean?}
\end{center}

\vspace{2em}

\begin{columns}[T]
\begin{column}{0.48\textwidth}
\textbf{Computational View:}
\begin{itemize}
    \item Vectors in high-dimensional space
    \item Learned from text co-occurrence
    \item Distributional patterns
    \item Statistical relationships
    \item "You shall know a word by the company it keeps"
\end{itemize}

\vspace{0.5em}
\alert{But is this how humans represent meaning?}
\end{column}

\begin{column}{0.48\textwidth}
\textbf{Human View:}
\begin{itemize}
    \item Sensory experiences
    \item Emotional associations
    \item Physical grounding
    \item Social context
    \item Multimodal integration
    \item Embodied understanding
\end{itemize}

\vspace{0.5em}
\alert{Can text-only models capture this?}
\end{column}
\end{columns}

\vspace{1em}

\begin{center}
\textit{This lecture explores the relationship between computational and cognitive semantics}
\end{center}
\end{frame}

% ============================================================
% HUMAN SEMANTICS
% ============================================================
\begin{frame}{How Do Humans Represent Meaning? 🧠}
\textbf{Example: The word "coffee"}

\vspace{1em}

\begin{columns}[T]
\begin{column}{0.48\textwidth}
\textbf{What comes to YOUR mind?}
\begin{itemize}
    \item \textbf{Visual:} brown liquid, mug, steam, beans
    \item \textbf{Olfactory:} aroma, roasted smell
    \item \textbf{Gustatory:} bitter taste, smooth texture
    \item \textbf{Tactile:} hot, warm cup, liquid
    \item \textbf{Auditory:} brewing sounds, pouring
    \item \textbf{Motor:} lifting cup, drinking motion
    \item \textbf{Contextual:} morning routine, work, café
    \item \textbf{Emotional:} comfort, alertness, pleasure
    \item \textbf{Social:} conversations, meetings
    \item \textbf{Cultural:} Starbucks, espresso, traditions
\end{itemize}
\end{column}

\begin{column}{0.48\textwidth}
\textbf{Word2Vec representation:}

\vspace{0.5em}
\texttt{coffee = [0.23, -0.45, 0.67, 0.12, -0.89, ...]}

\vspace{0.5em}
(300 numbers learned from text)

\vspace{1em}
\textbf{What's missing?}
\begin{itemize}
    \item No sensory grounding
    \item No embodied experience
    \item No emotional content
    \item No perceptual features
    \item No action affordances
    \item Purely linguistic
\end{itemize}

\vspace{1em}
\begin{alertblock}{The Grounding Problem}
Computational models lack the rich, multimodal, embodied understanding that humans have.
\end{alertblock}
\end{column}
\end{columns}
\end{frame}

% ============================================================
% SYMBOL GROUNDING
% ============================================================
\begin{frame}{The Symbol Grounding Problem 🔗}
\textbf{Harnad (1990): Can symbols have intrinsic meaning?}

\vspace{1em}

\begin{exampleblock}{The Chinese Room Argument (Searle, 1980)}
\begin{itemize}
    \item Person in room, doesn't understand Chinese
    \item Receives Chinese symbols, has rule book
    \item Follows rules to produce Chinese responses
    \item From outside: appears to understand Chinese
    \item But has no understanding of meaning!
\end{itemize}

\vspace{0.5em}
\textbf{Analogy:} Language models manipulate symbols without grounding
\end{exampleblock}

\vspace{1em}

\begin{columns}[T]
\begin{column}{0.48\textwidth}
\textbf{Symbol Systems:}
\begin{itemize}
    \item Symbols refer to other symbols
    \item Purely syntactic manipulation
    \item No connection to world
    \item "Ungrounded" semantics
\end{itemize}
\end{column}

\begin{column}{0.48\textwidth}
\textbf{Grounded Systems:}
\begin{itemize}
    \item Symbols connected to perception
    \item Linked to action
    \item Embodied in physical world
    \item "Intrinsic" semantics
\end{itemize}
\end{column}
\end{columns}

\vspace{0.5em}
\small
\textit{References: Searle (1980). "Minds, Brains, and Programs"; Harnad (1990). "The Symbol Grounding Problem"}
\end{frame}

% ============================================================
% EMBODIED COGNITION
% ============================================================
\begin{frame}{Embodied Cognition Theory 🏃}
\textbf{Meaning arises from bodily experience and sensorimotor interaction}

\vspace{1em}

\begin{columns}[T]
\begin{column}{0.48\textwidth}
\textbf{Key Principles:}
\begin{enumerate}
    \item \textbf{Embodiment:} Cognition shaped by body
    \item \textbf{Situatedness:} Meaning context-dependent
    \item \textbf{Enactivism:} Knowing through doing
    \item \textbf{Grounding:} Concepts tied to perception/action
\end{enumerate}

\vspace{0.5em}
\textbf{Evidence:}
\begin{itemize}
    \item Reading "kick": activates motor cortex
    \item Understanding "grasp": simulates action
    \item Processing emotions: activates body maps
    \item Conceptual metaphors based on physical experience
\end{itemize}
\end{column}

\begin{column}{0.48\textwidth}
\begin{center}
\begin{tikzpicture}[scale=0.8]
    % Brain
    \draw[thick, fill=pink!30] (0,3) ellipse (1.5 and 1);
    \node at (0, 3) {Brain};

    % Body
    \draw[thick, fill=blue!20] (0,0.5) ellipse (1 and 1.5);
    \node at (0, 0.5) {Body};

    % Environment
    \draw[thick, fill=green!20] (-3,-1) rectangle (3,5);
    \node at (0, 5.5) {Environment};

    % Arrows
    \draw[<->, very thick, red] (0, 2) -- (0, 1.7);
    \draw[<->, very thick, blue] (-0.8, 0.5) -- (-2, 2);
    \draw[<->, very thick, blue] (0.8, 0.5) -- (2, 2);

    % Labels
    \node at (2.5, 1.5) {Perception};
    \node at (-2.5, 1.5) {Action};
    \node at (1, 1.85) {Embodied};
\end{tikzpicture}
\end{center}

\vspace{0.5em}
\textbf{Implications for AI:}
\begin{itemize}
    \item Text-only models miss embodiment
    \item Need multimodal grounding
    \item Interaction with environment
    \item Sensorimotor foundations
\end{itemize}
\end{column}
\end{columns}

\vspace{0.5em}
\small
\textit{Reference: Barsalou (2008). "Grounded Cognition"}
\end{frame}

% ============================================================
% DISTRIBUTIONAL HYPOTHESIS
% ============================================================
\begin{frame}{The Distributional Hypothesis Revisited 📖}
\begin{center}
\huge
\textit{"You shall know a word by the company it keeps"}
\end{center}

\begin{center}
--- J.R. Firth (1957)
\end{center}

\vspace{1em}

\textbf{Strong version:} Word meaning IS distributional patterns

\textbf{Weak version:} Distributional patterns REFLECT meaning

\vspace{1em}

\begin{columns}[T]
\begin{column}{0.48\textwidth}
\textbf{Supports:}
\begin{itemize}
    \item Works remarkably well in practice
    \item Captures semantic similarity
    \item Enables analogical reasoning
    \item Scales to huge vocabularies
    \item Unsupervised learning
    \item Aligned with usage-based linguistics
\end{itemize}
\end{column}

\begin{column}{0.48\textwidth}
\textbf{Limitations:}
\begin{itemize}
    \item Correlation $\neq$ causation
    \item Lacks perceptual grounding
    \item No embodied understanding
    \item Missing common sense
    \item Can't handle novel situations
    \item Reflects training data biases
\end{itemize}
\end{column}
\end{columns}

\vspace{1em}

\begin{alertblock}{Key Question}
Is distributional semantics sufficient for meaning, or just a useful approximation?
\end{alertblock}

\small
\textit{Reference: Boleda (2020). "Distributional Semantics and Linguistic Theory"}
\end{frame}

% ============================================================
% SEMANTIC SIMILARITY
% ============================================================
\begin{frame}{What Does "Semantic Similarity" Really Mean? 🔍}
\textbf{Different types of similarity:}

\vspace{1em}

\begin{columns}[T]
\begin{column}{0.48\textwidth}
\textbf{1. Taxonomic:}
\begin{itemize}
    \item dog $\sim$ cat (both animals)
    \item red $\sim$ blue (both colors)
    \item "Is-a" relationships
\end{itemize}

\vspace{0.3em}
\textbf{2. Thematic:}
\begin{itemize}
    \item dog $\sim$ leash (co-occur in events)
    \item coffee $\sim$ cup (functional relation)
    \item "Goes-with" relationships
\end{itemize}

\vspace{0.3em}
\textbf{3. Synonymy:}
\begin{itemize}
    \item big $\sim$ large (same meaning)
    \item happy $\sim$ joyful
\end{itemize}

\vspace{0.3em}
\textbf{4. Functional:}
\begin{itemize}
    \item knife $\sim$ scissors (similar function)
    \item chair $\sim$ stool
\end{itemize}
\end{column}

\begin{column}{0.48\textwidth}
\textbf{What do models capture?}

\vspace{0.5em}
\begin{itemize}
    \item Word2Vec: Mostly thematic + taxonomic
    \item BERT: Better at taxonomic
    \item Varies by context window size
    \item Training data matters!
\end{itemize}

\vspace{1em}
\textbf{Human judgments:}
\begin{itemize}
    \item SimLex-999: True similarity
    \item WordSim-353: Relatedness (broader)
    \item Models better at relatedness than similarity
\end{itemize}

\vspace{1em}
\begin{exampleblock}{Example}
\textbf{Similar:} car $\sim$ automobile (0.9)

\textbf{Related:} car $\sim$ road (0.7)

\vspace{0.3em}
Models often conflate these!
\end{exampleblock}
\end{column}
\end{columns}

\vspace{0.5em}
\small
\textit{Reference: Hill et al. (2015). "SimLex-999: Evaluating Semantic Models with Genuine Similarity Estimation"}
\end{frame}

% ============================================================
% GRAND ET AL
% ============================================================
\begin{frame}{Semantic Projection: Recovering Human Knowledge 🔬}
\textbf{Grand et al. (2022): Can we extract human-like features from embeddings?}

\vspace{1em}

\begin{columns}[T]
\begin{column}{0.48\textwidth}
\textbf{The Experiment:}
\begin{enumerate}
    \item Collect human ratings on perceptual features
    \begin{itemize}
        \item Is it edible?
        \item Is it heavy?
        \item Is it alive?
        \item Can you hold it?
    \end{itemize}
    \item Train linear projections from word embeddings
    \item Test if embeddings predict human ratings
\end{enumerate}

\vspace{0.5em}
\textbf{Key Findings:}
\begin{itemize}
    \item Embeddings encode surprisingly rich knowledge
    \item Can predict perceptual features
    \item Even vision and motor properties!
    \item Better with larger models
\end{itemize}
\end{column}

\begin{column}{0.48\textwidth}
\begin{center}
\begin{tikzpicture}[scale=0.7]
    % Word Embedding
    \draw[fill=blue!20] (0,3) rectangle (4,4);
    \node at (2, 3.5) {Word Embedding};
    \node[font=\tiny] at (2, 3) {$[0.23, -0.45, ..., 0.12]$};

    % Projection
    \draw[->, very thick] (2, 3) -- (2, 2);
    \node at (3.5, 2.5) {Linear};
    \node at (3.5, 2.2) {Projection};

    % Features
    \draw[fill=green!20] (0,0.3) rectangle (4,1.5);
    \node[font=\small] at (2, 1.3) {Perceptual Features};
    \node[font=\tiny] at (2, 0.9) {Size: 0.7};
    \node[font=\tiny] at (2, 0.6) {Edible: 0.1};

    % Comparison
    \draw[->, thick] (4.2, 0.9) -- (5.5, 0.9);
    \draw[fill=yellow!20] (5.7,0.3) rectangle (8,1.5);
    \node[font=\tiny] at (6.85, 1) {Human};
    \node[font=\tiny] at (6.85, 0.6) {Ratings};

    % Correlation
    \node at (6.85, -0.3) {$r \approx 0.7$};
\end{tikzpicture}
\end{center}

\vspace{0.5em}
\textbf{Interpretation:}
\begin{itemize}
    \item Text co-occurrence captures real-world properties
    \item Indirect grounding through language
    \item But still not true perceptual grounding
\end{itemize}
\end{column}
\end{columns}

\vspace{0.5em}
\small
\textit{Reference: Grand et al. (2022). "Semantic projection recovers rich human knowledge of multiple object features from word embeddings"}
\end{frame}

% ============================================================
% MULTIMODAL
% ============================================================
\begin{frame}{Multimodal Models: Bridging the Gap 🌉}
\textbf{Combining language with perception}

\vspace{1em}

\begin{columns}[T]
\begin{column}{0.48\textwidth}
\textbf{Vision-Language Models:}
\begin{itemize}
    \item CLIP (OpenAI)
    \item ALIGN (Google)
    \item Flamingo (DeepMind)
    \item GPT-4V (OpenAI)
\end{itemize}

\vspace{0.5em}
\textbf{Key Idea:}
\begin{itemize}
    \item Learn joint embedding space
    \item Text and images map to same space
    \item Enables cross-modal understanding
    \item Grounds language in vision
\end{itemize}

\vspace{0.5em}
\textbf{Training:}
\begin{itemize}
    \item Image-caption pairs
    \item Contrastive learning
    \item Matching images to descriptions
    \item Large-scale (400M+ pairs)
\end{itemize}
\end{column}

\begin{column}{0.48\textwidth}
\begin{center}
\begin{tikzpicture}[scale=0.7]
    % Image
    \draw[fill=blue!20] (0,4) rectangle (2,6);
    \node at (1, 5) {Image};

    % Text
    \draw[fill=green!20] (4,4) rectangle (6,6);
    \node at (5, 5) {Text};

    % Encoders
    \draw[->, thick] (1, 4) -- (1, 3);
    \draw[->, thick] (5, 4) -- (5, 3);

    \node[font=\small] at (1, 3.5) {Vision};
    \node[font=\small] at (1, 3.2) {Encoder};

    \node[font=\small] at (5, 3.5) {Text};
    \node[font=\small] at (5, 3.2) {Encoder};

    % Joint Space
    \draw[fill=purple!20] (0,0.5) rectangle (6,2.5);
    \node at (3, 1.5) {Joint Embedding Space};

    \draw[->, thick] (1, 3) -- (2, 2.5);
    \draw[->, thick] (5, 3) -- (4, 2.5);

    % Points in space
    \fill[blue] (1.5, 1.8) circle (3pt);
    \fill[green] (1.7, 1.6) circle (3pt);
    \node[font=\tiny] at (1.6, 1.2) {Similar!};

    \fill[blue] (4.5, 1) circle (3pt);
    \fill[green] (3.8, 1.3) circle (3pt);
\end{tikzpicture}
\end{center}

\vspace{0.5em}
\textbf{Benefits:}
\begin{itemize}
    \item Visual grounding
    \item Zero-shot classification
    \item Better generalization
    \item More "human-like"
\end{itemize}
\end{column}
\end{columns}

\vspace{0.5em}
\small
\textit{Reference: Radford et al. (2021). "Learning Transferable Visual Models From Natural Language Supervision" (CLIP)}
\end{frame}

% ============================================================
% CONCEPTUAL SPACES
% ============================================================
\begin{frame}{Conceptual Spaces Theory 🌌}
\textbf{Gärdenfors (2000): Meaning as geometry}

\vspace{1em}

\begin{columns}[T]
\begin{column}{0.48\textwidth}
\textbf{Key Ideas:}
\begin{itemize}
    \item Concepts represented in quality dimensions
    \item Dimensions: color, size, temperature, etc.
    \item Each dimension has a metric
    \item Concepts are regions in space
    \item Similarity = geometric proximity
\end{itemize}

\vspace{0.5em}
\textbf{Example: Colors}
\begin{itemize}
    \item Hue, saturation, brightness
    \item Natural categories (red, blue, green)
    \item Fuzzy boundaries
    \item Prototypes at centers
\end{itemize}
\end{column}

\begin{column}{0.48\textwidth}
\begin{center}
\begin{tikzpicture}[scale=0.8]
    % Axes
    \draw[->] (0,0) -- (5,0) node[right] {Size};
    \draw[->] (0,0) -- (0,4) node[above] {Danger};

    % Regions
    \draw[fill=red!20, opacity=0.5] (0.5,3) circle (0.8);
    \node at (0.5, 3) {lion};

    \draw[fill=blue!20, opacity=0.5] (1,1) circle (0.6);
    \node at (1, 1) {cat};

    \draw[fill=green!20, opacity=0.5] (4,3) circle (0.9);
    \node at (4, 3) {elephant};

    \draw[fill=yellow!20, opacity=0.5] (3.5,0.8) circle (0.5);
    \node[font=\small] at (3.5, 0.8) {rabbit};

    % Prototype
    \fill[red] (0.5, 3) circle (2pt);
    \fill[blue] (1, 1) circle (2pt);
\end{tikzpicture}
\end{center}

\vspace{0.5em}
\textbf{Relation to Embeddings:}
\begin{itemize}
    \item Similar geometric structure
    \item But learned dimensions
    \item Not interpretable qualities
    \item No grounding in perception
\end{itemize}
\end{column}
\end{columns}

\vspace{0.5em}
\small
\textit{Reference: Gärdenfors (2000). "Conceptual Spaces: The Geometry of Thought"}
\end{frame}

% ============================================================
% LEXICAL SEMANTICS
% ============================================================
\begin{frame}{Lexical Semantic Theories 📚}
\textbf{How do linguists think about word meaning?}

\vspace{1em}

\begin{columns}[T]
\begin{column}{0.48\textwidth}
\textbf{1. Feature-Based:}
\begin{itemize}
    \item Words = bundles of features
    \item bachelor = [+human, +male, +adult, -married]
    \item Compositional
    \item Logical
\end{itemize}

\vspace{0.3em}
\textbf{2. Prototype Theory:}
\begin{itemize}
    \item Categories have best examples
    \item Robin is a prototypical bird
    \item Penguin is peripheral
    \item Graded membership
\end{itemize}

\vspace{0.3em}
\textbf{3. Frame Semantics:}
\begin{itemize}
    \item Words evoke conceptual frames
    \item "Buy" activates commerce frame
    \item Buyer, seller, goods, money
    \item FrameNet resource
\end{itemize}
\end{column}

\begin{column}{0.48\textwidth}
\textbf{4. Construction Grammar:}
\begin{itemize}
    \item Meaning from form-function pairings
    \item Idioms, patterns
    \item Usage-based
\end{itemize}

\vspace{0.3em}
\textbf{5. Word Sense Disambiguation:}
\begin{itemize}
    \item Words have multiple senses
    \item WordNet: hierarchical ontology
    \item Bank$_1$ (financial), Bank$_2$ (riverside)
\end{itemize}

\vspace{1em}
\begin{alertblock}{Question}
Which theories align with distributional models?

\vspace{0.3em}
Answer: Mostly usage-based views (Construction Grammar, Prototype Theory)

\vspace{0.3em}
Struggle with: Feature analysis, Frame semantics
\end{alertblock}
\end{column}
\end{columns}

\vspace{0.5em}
\small
\textit{References: Fillmore (1982). "Frame Semantics"; Rosch (1975). "Cognitive Representations of Semantic Categories"}
\end{frame}

% ============================================================
% EMPIRICAL EVIDENCE
% ============================================================
\begin{frame}{Empirical Evidence from Neuroscience 🧬}
\textbf{What does the brain tell us about semantic representation?}

\vspace{1em}

\begin{columns}[T]
\begin{column}{0.48\textwidth}
\textbf{fMRI Studies:}
\begin{itemize}
    \item Can predict brain activity from word embeddings
    \item Semantic information distributed across cortex
    \item Different regions for different features
    \item Temporal lobe: objects
    \item Motor cortex: actions
    \item Visual cortex: visual features
\end{itemize}

\vspace{0.5em}
\textbf{Findings:}
\begin{itemize}
    \item Word2Vec correlates with neural patterns
    \item But: imperfect match
    \item Human brain uses multimodal integration
    \item Language areas + sensory areas
\end{itemize}
\end{column}

\begin{column}{0.48\textwidth}
\textbf{Brain vs. Model Representations:}

\vspace{0.5em}
\begin{tabular}{lcc}
\toprule
\textbf{Feature} & \textbf{Brain} & \textbf{Model} \\
\midrule
Multimodal & ✓ & ✗ \\
Distributed & ✓ & ✓ \\
Hierarchical & ✓ & ✓ \\
Context-sensitive & ✓ & ✓ (BERT) \\
Grounded & ✓ & ✗ \\
Fast & ✓ & ✓ \\
\bottomrule
\end{tabular}

\vspace{1em}
\textbf{Key Insight:}
\begin{itemize}
    \item Models capture some aspects
    \item But miss crucial grounding
    \item Convergent representation learning?
    \item Or fundamentally different?
\end{itemize}
\end{column}
\end{columns}

\vspace{0.5em}
\small
\textit{References: Mitchell et al. (2008). "Predicting Human Brain Activity Associated with the Meanings of Nouns"; Huth et al. (2016). "Natural speech reveals the semantic maps that tile human cerebral cortex"}
\end{frame}

% ============================================================
% STOCHASTIC PARROTS
% ============================================================
\begin{frame}{The "Stochastic Parrots" Debate 🦜}
\textbf{Bender et al. (2021): On the Dangers of Stochastic Parrots}

\vspace{1em}

\begin{columns}[T]
\begin{column}{0.48\textwidth}
\textbf{The Argument:}
\begin{itemize}
    \item LLMs learn form, not meaning
    \item "Stochastic parrots" - repeating patterns
    \item No understanding of world
    \item No communicative intent
    \item Risk: Mistaking fluency for understanding
\end{itemize}

\vspace{0.5em}
\textbf{Evidence:}
\begin{itemize}
    \item Fail on simple reasoning
    \item Sensitive to phrasing
    \item Hallucinate facts
    \item No common sense
    \item Brittle to adversarial inputs
\end{itemize}
\end{column}

\begin{column}{0.48\textwidth}
\textbf{Counter-Arguments:}
\begin{itemize}
    \item Emergent capabilities at scale
    \item Performance on complex tasks
    \item Transfer to new domains
    \item Maybe understanding = prediction?
    \item Pragmatic success criterion
\end{itemize}

\vspace{1em}
\begin{alertblock}{The Core Question}
\textbf{Can meaning arise from form alone?}

\vspace{0.3em}
Or do we need grounding in:
\begin{itemize}
    \item Perception?
    \item Action?
    \item Social interaction?
    \item Physical embodiment?
\end{itemize}
\end{alertblock}
\end{column}
\end{columns}

\vspace{0.5em}
\small
\textit{References: Bender \& Koller (2020). "Climbing towards NLU"; Bender et al. (2021). "On the Dangers of Stochastic Parrots"}
\end{frame}

% ============================================================
% COMMON SENSE
% ============================================================
\begin{frame}{Common Sense Reasoning 🤯}
\textbf{What humans know but models don't}

\vspace{1em}

\begin{columns}[T]
\begin{column}{0.48\textwidth}
\textbf{Examples:}

\vspace{0.3em}
\begin{exampleblock}{Physical}
\begin{itemize}
    \item Water is wet
    \item Objects fall down
    \item You can't be in two places at once
    \item Heavy things are hard to lift
\end{itemize}
\end{exampleblock}

\vspace{0.3em}
\begin{exampleblock}{Social}
\begin{itemize}
    \item People have feelings
    \item Insults are hurtful
    \item Promises should be kept
    \item Eye contact shows attention
\end{itemize}
\end{exampleblock}
\end{column}

\begin{column}{0.48\textwidth}
\textbf{Why Models Struggle:}
\begin{itemize}
    \item Not explicit in text
    \item Assumed background knowledge
    \item Requires world experience
    \item Needs causal reasoning
    \item Embodied understanding
\end{itemize}

\vspace{1em}
\textbf{Benchmark Datasets:}
\begin{itemize}
    \item Winograd Schema Challenge
    \item PIQA (Physical Interaction QA)
    \item SocialIQA
    \item CommonsenseQA
\end{itemize}

\vspace{0.5em}
\textbf{Progress:}
\begin{itemize}
    \item Large models do better
    \item But still far from human
    \item Often exploit shortcuts
    \item Memorization vs. reasoning?
\end{itemize}
\end{column}
\end{columns}
\end{frame}

% ============================================================
% COMPOSITIONALITY
% ============================================================
\begin{frame}{Compositionality: Phrases and Sentences 🧩}
\textbf{How do we combine word meanings?}

\vspace{1em}

\begin{columns}[T]
\begin{column}{0.48\textwidth}
\textbf{The Problem:}
\begin{itemize}
    \item "hot" + "dog" $\neq$ "hot dog"
    \item "red" + "herring" $\neq$ "red herring"
    \item Idioms, metaphors, collocations
    \item Non-compositional meaning
\end{itemize}

\vspace{0.5em}
\textbf{Classical Approaches:}
\begin{itemize}
    \item Vector addition: $\vec{hot} + \vec{dog}$
    \item Element-wise multiplication
    \item Tensor products
    \item Fails for non-compositional phrases!
\end{itemize}

\vspace{0.5em}
\textbf{Neural Approaches:}
\begin{itemize}
    \item RNNs, LSTMs
    \item Transformers (BERT)
    \item Learn composition functions
    \item Better but not perfect
\end{itemize}
\end{column}

\begin{column}{0.48\textwidth}
\textbf{Compositionality Types:}

\vspace{0.3em}
\textbf{1. Semantic Composition:}
\begin{itemize}
    \item "big red ball" = big $\cap$ red $\cap$ ball
    \item Intersective
\end{itemize}

\vspace{0.3em}
\textbf{2. Non-intersective:}
\begin{itemize}
    \item "fake gun" $\neq$ fake $\cap$ gun
    \item Adjective changes interpretation
\end{itemize}

\vspace{0.3em}
\textbf{3. Metaphorical:}
\begin{itemize}
    \item "Time is money"
    \item Cross-domain mapping
\end{itemize}

\vspace{1em}
\begin{alertblock}{Current Status}
Transformers handle many cases through attention, but still struggle with:
\begin{itemize}
    \item Novel compositions
    \item Systematic generalization
    \item Logical operators (negation!)
\end{itemize}
\end{alertblock}
\end{column}
\end{columns}
\end{frame}

% ============================================================
% DISCUSSION
% ============================================================
\begin{frame}{Grand Discussion 💬}
\begin{center}
\Large
\textbf{Do large language models "understand" language?}
\end{center}

\vspace{1em}

\begin{columns}[T]
\begin{column}{0.48\textwidth}
\textbf{Arguments FOR:}
\begin{itemize}
    \item Solve complex tasks
    \item Generalize to new domains
    \item Show emergent capabilities
    \item Capture linguistic structure
    \item Pragmatic criterion: if it works...
    \item Maybe understanding = prediction
    \item Human understanding also imperfect
\end{itemize}

\vspace{0.5em}
\textit{"The question is not whether machines think, but whether they behave intelligently" - Turing}
\end{column}

\begin{column}{0.48\textwidth}
\textbf{Arguments AGAINST:}
\begin{itemize}
    \item No grounding in reality
    \item No embodied experience
    \item No intentionality
    \item Brittle, exploit shortcuts
    \item Hallucinate confidently
    \item No causal reasoning
    \item Missing common sense
    \item Form without meaning
\end{itemize}

\vspace{0.5em}
\textit{"Understanding requires grounding in perception and action" - Embodied cognition}
\end{column}
\end{columns}

\vspace{1em}

\begin{alertblock}{Perhaps the wrong question?}
Instead of "Do they understand?", ask:
\begin{itemize}
    \item What do they represent?
    \item How does it differ from humans?
    \item What are the limits?
    \item How can we improve?
\end{itemize}
\end{alertblock}
\end{frame}

% ============================================================
% FUTURE DIRECTIONS
% ============================================================
\begin{frame}{Future Directions 🔮}
\textbf{Bridging computational and cognitive semantics}

\vspace{1em}

\begin{enumerate}
    \item \textbf{Multimodal Learning:}
    \begin{itemize}
        \item Vision + language (CLIP, Flamingo)
        \item Audio, tactile, proprioception
        \item Grounding in multiple modalities
    \end{itemize}

    \item \textbf{Embodied AI:}
    \begin{itemize}
        \item Robots learning through interaction
        \item Simulation environments
        \item Sensorimotor grounding
    \end{itemize}

    \item \textbf{Neurosymbolic AI:}
    \begin{itemize}
        \item Combining neural and symbolic
        \item Logical reasoning + learning
        \item Structured knowledge graphs
    \end{itemize}

    \item \textbf{Cognitively-Inspired Architectures:}
    \begin{itemize}
        \item Memory systems
        \item Attention mechanisms
        \item Episodic learning
        \item Meta-learning
    \end{itemize}

    \item \textbf{Social Grounding:}
    \begin{itemize}
        \item Learning through dialogue
        \item Cultural context
        \item Pragmatic understanding
    \end{itemize}
\end{enumerate}
\end{frame}

% ============================================================
% PRACTICAL IMPLICATIONS
% ============================================================
\begin{frame}{Practical Implications 💼}
\textbf{What this means for NLP practitioners}

\vspace{1em}

\begin{enumerate}
    \item \textbf{Know the Limits:}
    \begin{itemize}
        \item Models excel at pattern matching
        \item Struggle with true reasoning
        \item Need task-specific evaluation
        \item Don't assume understanding
    \end{itemize}

    \item \textbf{Choose Appropriate Tasks:}
    \begin{itemize}
        \item Good: classification, retrieval, similarity
        \item Moderate: summarization, translation
        \item Challenging: reasoning, planning, common sense
    \end{itemize}

    \item \textbf{Augment with Structure:}
    \begin{itemize}
        \item Knowledge graphs
        \item Rules and constraints
        \item Domain expertise
        \item Human-in-the-loop
    \end{itemize}

    \item \textbf{Evaluate Carefully:}
    \begin{itemize}
        \item Beyond accuracy metrics
        \item Test edge cases
        \item Adversarial robustness
        \item Bias and fairness
        \item Interpretability
    \end{itemize}

    \item \textbf{Stay Informed:}
    \begin{itemize}
        \item Rapidly evolving field
        \item New capabilities emerging
        \item Ethical considerations
        \item Interdisciplinary insights
    \end{itemize}
\end{enumerate}
\end{frame}

% ============================================================
% SUMMARY
% ============================================================
\begin{frame}{Summary 🎯}
\textbf{What we learned today:}

\begin{enumerate}
    \item \textbf{Symbol Grounding:} Computational models lack perceptual grounding

    \item \textbf{Embodied Cognition:} Human meaning tied to bodily experience

    \item \textbf{Distributional Semantics:} Powerful but incomplete theory

    \item \textbf{Semantic Similarity:} Multiple types, models capture some

    \item \textbf{Empirical Evidence:} Models align with neural patterns but miss multimodality

    \item \textbf{Current Debates:}
    \begin{itemize}
        \item Stochastic parrots vs. emergent understanding
        \item Form vs. meaning
        \item Prediction vs. comprehension
    \end{itemize}

    \item \textbf{Future:} Multimodal, embodied, neurosymbolic approaches
\end{enumerate}

\vspace{1em}

\begin{center}
\Large
\textbf{The gap between computation and cognition remains,}

\textbf{but we're making progress! 🌉}
\end{center}
\end{frame}

% ============================================================
% REFERENCES
% ============================================================
\begin{frame}{Key References 📚}
\small
\textbf{Foundational Papers:}
\begin{itemize}
    \item Searle, J. (1980). "Minds, Brains, and Programs"
    \item Harnad, S. (1990). "The Symbol Grounding Problem"
    \item Barsalou, L. (2008). "Grounded Cognition"
    \item Gärdenfors, P. (2000). "Conceptual Spaces: The Geometry of Thought"
\end{itemize}

\vspace{0.3em}
\textbf{Distributional Semantics:}
\begin{itemize}
    \item Firth, J.R. (1957). "A Synopsis of Linguistic Theory"
    \item Boleda, G. (2020). "Distributional Semantics and Linguistic Theory"
    \item Hill et al. (2015). "SimLex-999"
    \item Grand et al. (2022). "Semantic Projection Recovers Rich Human Knowledge"
\end{itemize}

\vspace{0.3em}
\textbf{Critical Perspectives:}
\begin{itemize}
    \item Bender \& Koller (2020). "Climbing towards NLU: On Meaning, Form, and Understanding"
    \item Bender et al. (2021). "On the Dangers of Stochastic Parrots"
\end{itemize}

\vspace{0.3em}
\textbf{Neuroscience:}
\begin{itemize}
    \item Mitchell et al. (2008). "Predicting Human Brain Activity"
    \item Huth et al. (2016). "Natural Speech Reveals Semantic Maps"
\end{itemize}

\vspace{0.3em}
\textbf{Multimodal:}
\begin{itemize}
    \item Radford et al. (2021). "Learning Transferable Visual Models" (CLIP)
\end{itemize}
\end{frame}

% ============================================================
% QUESTIONS
% ============================================================
\begin{frame}{Final Thoughts 🌟}
\begin{center}
\LARGE
\textbf{As we build ever-larger language models,}

\vspace{0.5em}
\textbf{let's not forget to ask:}

\vspace{1em}
\textit{What are they really learning?}

\vspace{0.5em}
\textit{How does it compare to human understanding?}

\vspace{0.5em}
\textit{What's missing?}

\vspace{0.5em}
\textit{How can we do better?}

\vspace{2em}
\textbf{Questions? 🙋}
\end{center}
\end{frame}

\end{document}
